\section{Introduction}\label{sec:intro119}
Dynamic investment under risk requires a continuation value, a feasible action,
and a decision about how much information to acquire before acting. A small
continuation fitting error need not imply a small decision loss. It can move a
nearly binding action, alter the adverse innovation tail, or leave an expensive
original-model verification problem unresolved. Neural Bellman Operators (NBOs)
study these objects together: recursively constructed and learned continuation
maps propose policies, their own future decisions determine policy value, and
original-law Bellman inequalities qualify the actions actually delivered.

The economic question is when an acquired neural continuation helps deliver
original-objective decisions at a competitive complete cost. Answering it
requires conventional representations with the same information, feasible
actions and numerical components. A conventional hinge encoding of the same
ReLU map receives the same policy and certificate. A quadratic continuation
that wins a complete-work comparison remains a successful control. The purpose
of these comparisons is to identify which part of continuation construction,
information acquisition and policy improvement creates economic value.

This paper develops a chain from recursive continuation construction to
action-aligned information and verified decisions. The retained own-future NBO
construction gives an exactly realizable neural continuation and a measurable
feasible actor without assuming successful nonconvex training. Centered
recursive-risk inequalities separate irrelevant value-level shifts from
action-dependent approximation errors. Original-model supporting quadratics
then connect the realized constrained action to the original Bellman optimum.
The assumptions behind Bellman existence, measurability, continuity and
uniqueness are stated explicitly; the supplement proves them for compact
innovation supports and verifies the original investment primitives.

Nonlinear risk introduces an economic obstacle beyond directional fitting
error. Even when two continuations have identical action derivatives at every
innovation, their maximizing risk densities can differ. The retained
continuous-law tail-switch example has zero directional approximation error
and original action loss $1/4$. Centered quantile-band bounds identify the
additional near-tail information needed to control this discrepancy. The
present addition determines how much point information is needed to resolve
that tail, and constructs a policy certificate when its location is only
approximately known.

For a convex $K$-Lipschitz continuation on $[-R,R]$ under the original uniform
law, an adverse tail of probability $t=1-\alpha$ consists of at most two
endpoint intervals. Its complement is a central interval of fixed length.
The derivative of the integral over that central interval is the difference
between two endpoint continuation values; it is monotone and Lipschitz.
Consequently paid bisection locates the relevant tails, and midpoint--secant
integration can concentrate on their total length $2Rt$. Including every
localization query, we prove a deterministic worst-case risk-width upper bound
of order $KRtN^{-2}$. Together with a fully restated indistinguishability
lower bound, this gives matching risk-level dependence for $N\ge32$. This
resolves the factor $t^{-1}$ in the retained whole-support upper account.
The theorem specifies exact point information; bit precision and the work
needed to obtain each continuation value have their own implementation account.

The original nonlinear investment economy supplies an implementable version.
Primitive curvature implies that the optimal continuation along its innovation
direction is strongly convex with modulus $181/64$. For any proposed feasible
tail cut, outward intervals for the two endpoint values bound its failure to
maximize the original risk functional by a constrained quadratic calculation.
This defect is paid explicitly. The same fixed feasible density and global
continuation supports produce a lower bound over the entire original action
interval. Selected-tail secants and the unchanged implemented-future allowance
produce an upper bound on the actual complete policy. No derivative through
a rounded actor, fitted sign assertion or hypothetical future optimizer obtains
certificate authority.

The computational contribution is therefore a change in which original
continuation values are purchased and how their errors enter a decision,
rather than an internal fit-score improvement. A geometric precision schedule
can stop at the first qualified resolution, while a complete fallback handles
an intrinsically inadequate candidate or supporting minorant. The retained
completed-incumbent rule gives strict original-policy improvement. Its
remaining-gain budget bounds the economic improvement sacrificed by stopping
further search. A coupled supplementary account integrates common-density
policy comparisons, atom-safe tail remainders and signed recursive gain
transport, with the independently frozen earlier research on that question
credited in the development companion.

The numerical study holds the original two-date investment technology,
continuous innovation distribution, nested AVaR preference, feasible actions,
terminal controller and accuracy targets fixed. Eight arms share fresh paid
terminal-value labels: the online semi-analytic Bellman control, a rank-six
quadratic, ReLU, its equivalent conventional hinge, coordinate splines, full
quadratic and cubic polynomials, and tensor-grid interpolation. Each arm
receives the same candidate policies under each contemporaneous certificate
service. Three separately frozen phases retain 32,256 primary decisions and
50,176 fresh-process timing decisions, together with acquisition, fitting
audits, rejected attempts, repairs, original-law reference calculations and
failures. Phase A establishes the active-tail implementation; its adverse
elapsed-work comparison motivates phase B's bounded preliminary acquisition.
Phase B's remaining overhead motivates phase C's packed original evaluation.
The latter two phases are explicit development on the previously observed
catalogue. Every control is executed afresh, and neither phase replaces the
observations of its predecessor.

Phase C's packed service is faster than its uncombined budgeted control in
all 16 median-clock cells and faster than the retained geometric service in
ten, with all six losses reported. The ReLU service-time ratios against the
geometric control are $0.8614$ and $0.8494$ at the two risk levels. Returned
policies are unchanged across services, so their economic cost contrasts are
exactly zero. With fresh standalone acquisition and actual model loading
charged, the quadratic control still has the lower original mean economic
cost and lower complete component cost at both risk levels. Conventional
hinge and ReLU maps remain policy-identical. The fixed catalogue is a
reproducible numerical comparison; its results carry no population-sampling
interpretation. Seven fresh timing processes characterize observed machine
variation and do not supply sampling confidence intervals.

The contribution builds on established risk-averse dynamic programming and
policy iteration \citep{ruszczynski2010risk,ruszczynski2014erratum}, the AVaR
variational formulation \citep{rockafellaruryasev2002}, and information-based
integration \citep{Novak2015Integration97}. Inexact-oracle optimization and
decision-focused learning clarify related accuracy and acquisition questions
\citep{DevolderGlineurNesterov2014,ElmachtoubGrigas2022}. Here the mathematical
addition is the joint tail-location and integration account with matching
tail-length dependence, and its outward implementation as an all-action
certificate under the same original economic law. The compiler and these
certificates are available to every representation. Formula information and
derivative access are charged and distinguished from exact point values.

The original NBO program also concerns stochastic control and nonlinear
recursive utility. A Brownian Markov BSDE verification theorem states an
admissible connection, including driver measurability, the critic graph
$z=\sigma^{\mathsf T}\nabla u$, actor-chain derivatives, terminal data and
stopping-boundary errors. General NDU/FBSDE systems require their own existence,
information and objective assumptions. The investment evidence is attached to
its nested AVaR objective, while the recursive-utility theorem is attached to
the utility problem it actually proves.

The complete article retains earlier original-policy improvements and their
economic stopping costs, the historical quadratic complete-work winner, all
36 improved residual spans accompanied by an increase from 644 to 757
first-proposal rejections, and the unestablished neural-adoption regions.
These observations motivate action-aware acquisition and stronger numerical
methods. Current progress is evaluated by new proofs, complete decisions and
paid work. Neural-specific representation advantages, prospective acquisition
validation, longer horizons and calibrated economic adoption remain distinct
research obligations. The supplement supplies complete mathematical arguments;
the development companion and source archive preserve earlier theories,
applications, adverse evidence and independently frozen experiments.
